\section{What the i-RevNet reference implementation actually computes}
\label{app:irevnet}

\Cref{sec:flow} proposes using a pretrained, frozen i-RevNet \cite{jacobsen2018} as BB, relying on it being
an exactly invertible, dimension-preserving map with $\log|\det J_{\mathrm{BB}}(x)|=0$. This appendix records,
in the same spirit as \cref{app:code} for the Ng et al.\ reference implementation, what the released
i-RevNet code (\texttt{iRevNet.py}, \texttt{model\_utils.py}, \texttt{ILSVRC\_main.py}) actually computes,
using the authors' own terminology, and derives the Jacobian claim directly from that code rather than from
the architecture family in the abstract.

\subsection{Terminology}
\label{app:irevnet:terms}
The following terms are the authors' own, taken verbatim from the argument help-strings in
\texttt{ILSVRC\_main.py}:
\begin{quote}
\begin{flushleft}
\begin{tabular}{@{}l@{\quad}p{0.75\linewidth}@{}}
\texttt{--nBlocks}  & \emph{`number of blocks per stage'}\\
\texttt{--nStrides} & \emph{`number of strides per stage'}\\
\texttt{--nChannels}& \emph{`number of channels per stage'}\\
\texttt{--init\_ds} & \emph{`initial downsampling'}\\
\texttt{psi} & $\mathcal S_j$, concretely: it's the operator that does ``trade spatial size for channel 
count, losslessly'' -- same job as \texttt{init\_ds}'s squeeze, reused inside every stride-2 block so 
both streams' shapes keep matching $F$'s output after $F$ shrinks the spatial size. It's identity in the 
stride-1 blocks.
\end{tabular}
\end{flushleft}
\end{quote}
\begin{itemize}[leftmargin=1.5em]
\item \textbf{Stage.} One position in the \texttt{nBlocks}/\texttt{nStrides}/\texttt{nChannels} lists.
	  Each stage is a run of blocks at one fixed channel width, with only the first block doing work on 
	  shape.
      \texttt{iRevNet.irevnet\_stack} loops \texttt{zip(nChannels, nBlocks, nStrides)}, one
      (channels, depth, stride) triple per stage, expanded into \texttt{depth} blocks each.
\item \textbf{Block.} One \texttt{irevnet\_block} instance -- the actual coupling layer, in the sense of
      \cref{sec:flow:coupling}'s \cref{eq:additive-coupling}. Nothing coarser than this runs as a single
      computational step.
\item \textbf{Stride.} The help-string (`number of strides per stage') is a little misleading: it is not a
      count, but the one stride value used by a stage. Within a stage, only the \emph{first} block receives
      it (\texttt{strides = [stride] + [1]*(depth-1)} in \texttt{irevnet\_stack}); every other block in the
      stage is hardcoded to stride $1$.
\item \textbf{Channel} (an \texttt{nChannels} entry). The target channel width, per stream (the network
      always carries a pair of streams, \cref{app:irevnet:walkthrough}), for every block in that stage.
\item \textbf{Initial downsampling} (\texttt{init\_ds}). A one-time squeeze applied once, before stage $1$,
      via \texttt{self.init\_psi = psi(init\_ds)}; not itself a stage or a block.
\end{itemize}
\newpage

\subsection{The pretrained checkpoint actually used}
\label{app:irevnet:checkpoint}
The config distributed with the released, pretrained ImageNet checkpoint (the file accompanying
\url{https://drive.google.com/uc?id=1eQtZIpmtVM0-JDv_2nx3xMuGKjA1beLF&export=download} on
\cite{jacobsen2018}'s repository, there named \emph{i-RevNet 301}) is
\begin{quote}
\begin{flushleft}
\begin{quote}
\begin{flushleft}
\texttt{RevNet(nBlocks=[6, 16, 72, 6], nStrides=[2, 2, 2, 2],}\\
\texttt{\phantom{RevNet(}nChannels=[24, 96, 384, 1536],}\\
\texttt{\phantom{RevNet(}nClasses=1000, init\_ds=2, dropout\_rate=0.,}\\
\texttt{\phantom{RevNet(}affineBN=True, in\_shape=[3, 224, 224])}
\end{flushleft}
\end{quote}
\end{flushleft}
\end{quote}
i.e.\ $4$ stages, $100$ total blocks ($6+16+72+6$), reaching top-$1$/top-$5$ accuracy $74.018\%/91.590\%$.
This is the exact command \texttt{README.md} labels ``Train bijective i-RevNet on Imagenet,'' so the
checkpoint is the bijective (dimension-preserving), not the channel-widening injective, variant, despite
being considerably larger than either configuration printed in \cite{jacobsen2018}'s own Table~1 (a $29$M
bijective and a $181$M injective model). An OpenReview discussion thread on the paper independently confirms
the released \texttt{state\_dict} has approximately $119$M parameters -- consistent with this $100$-block
config and inconsistent with the $29$M Table~1 figure -- and separately notes that the released
\texttt{iRevNet.py} uses $3\times3$ kernels throughout each block's bottleneck, where the paper's Section~3.2
describes $1\times1,3\times3,1\times1$; the authors' reply states the released weights are ``a slightly
different version of the network described in the paper for internal reasons'' with unchanged conclusions,
without specifying what else differs. Two things are independently verifiable directly from the code,
however, regardless of that reply: (i) the model is dimension-preserving, checked in
\cref{app:irevnet:walkthrough}; (ii) it is genuinely bijective and not the padded/injective variant, checked
next.

\textbf{Bijective, not injective, confirmed from the channel schedule.} \texttt{irevnet\_block.\_\_init\_\_}
sets \texttt{self.pad = 2*out\_ch - in\_ch}, and only invokes \texttt{injective\_pad} (zero-padding extra
channels) when \texttt{self.pad != 0 and stride == 1}. For a stage's first block (\texttt{stride=2}), this
branch never fires regardless of \texttt{pad}'s value, since \texttt{psi} handles the shape change instead
(\cref{app:irevnet:walkthrough}). For every later, \texttt{stride=1} block in the same stage,
\texttt{irevnet\_stack} sets \texttt{in\_ch = 2 * channel} after each block, so
\texttt{pad = 2*out\_ch - in\_ch = 2*channel - 2*channel = 0} exactly -- \texttt{injective\_pad} never fires
anywhere in this config. The derivation of \cref{app:irevnet:jacobian} assumes this; a different
\texttt{nChannels} schedule that does trigger \texttt{injective\_pad} would need that step re-examined
separately, since it is genuinely dimension-expanding.
\newpage
\subsection{Forward-pass walkthrough}
\label{app:irevnet:walkthrough}
Traced against \texttt{iRevNet.forward} and \texttt{irevnet\_block.forward}, for the checkpoint of
\cref{app:irevnet:checkpoint} and a $3\times224\times224$ input. Everything below carries a \emph{pair} of
streams $(x_1,x_2)$ throughout; this is the coupling split of \cref{sec:flow:coupling}, nothing more.

\paragraph{Step 0: initial downsampling, before any stage.} \texttt{x = self.init\_psi.forward(x)}: 
\texttt{psi(init\_ds=2)} squeezes $3\times224\times224\to12\times112\times112$ (space$\to$channels, factor
$4$), then the result is split in half: $x_1,x_2$ at $6\times112\times112$ each.

\paragraph{Step 1: the four stages.} Only a stage's first block changes shape.
\begin{center}
\begin{tabular}{@{}lclcl@{}}
\toprule
stage & blocks & stride (1st block) & channels/stream & shape in $\to$ out (per stream)\\
\midrule
1 & 6  & 2, then $1\times5$  & 24   & $6\times112\times112\to24\times56\times56$\\
2 & 16 & 2, then $1\times15$ & 96   & $24\times56\times56\to96\times28\times28$\\
3 & 72 & 2, then $1\times71$ & 384  & $96\times28\times28\to384\times14\times14$\\
4 & 6  & 2, then $1\times5$  & 1536 & $384\times14\times14\to1536\times7\times7$\\
\bottomrule
\end{tabular}
\end{center}

\paragraph{Step 2: one block, stride $2$ (a stage's first block).}
\begin{quote}
\begin{flushleft}
\texttt{Fx2 = self.bottleneck\_block(x2)}\\
\texttt{if self.stride == 2:}\\
\texttt{\phantom{if }x1 = self.psi.forward(x1)}\\
\texttt{\phantom{if }x2 = self.psi.forward(x2)}\\
\texttt{y1 = Fx2 + x1}\\
\texttt{return (x2, y1)}
\end{flushleft}
\end{quote}
$F$'s own leading conv has \texttt{stride=stride} built in, so $F$ itself shrinks the spatial size,
computed on the \emph{pre}-\texttt{psi} $x_2$. \texttt{psi} is then applied to \emph{both} streams
afterward, purely to reshape them (space$\to$channels) to match $F$'s new, smaller output shape before the
addition; the table above confirms the reshaped $x_1,x_2$ and $Fx2$ land on the same shape at every stage.

\paragraph{Step 2$'$: one block, stride $1$ (every other block -- $96$ of the $100$ total).} \texttt{psi}
never runs; shapes never change. This is exactly $y_1=x_1,\ y_2=x_2+F(x_1)$ (up to which stream is labelled
$1$ vs.\ $2$), plain additive coupling as in \cref{eq:additive-coupling}, with no extra step.

\paragraph{The \texttt{first} flag.} Only the very first block overall (stage $1$, block $1$) skips its
\texttt{bottleneck\_block}'s leading BatchNorm+ReLU, since the input has just been split from raw pixels and
was never normalized by a prior block.

\paragraph{Kernel sizes.} Every convolution inside \texttt{bottleneck\_block} uses \texttt{kernel\_size=3}
(three $3\times3$ convolutions per block, bottleneck ratio \texttt{mult=4}); per
\cref{app:irevnet:checkpoint} this is not the $1\times1,3\times3,1\times1$ pattern printed in
\cite{jacobsen2018}'s own Section~3.2.

\paragraph{End of the stack.}
\begin{quote}
\begin{flushleft}
\texttt{out\_bij = merge(out[0], out[1])}\ \ \# concat -- $3072\times7\times7$\\
\texttt{out = F.relu(self.bn1(out\_bij))}\ \ \# classifier path only, from here\\
\texttt{out = F.avg\_pool2d(out, self.ds)}\\
\texttt{out = self.linear(out)}\\
\texttt{return out, out\_bij}
\end{flushleft}
\end{quote}
Two separate outputs. \texttt{out} is the $1000$-way ImageNet classifier logits -- not invertible, and not
what \cref{sec:flow} needs. \texttt{out\_bij} is the actual bijective representation, $3072\times7\times7$;
this, not \texttt{out}, is $\tilde x=\mathrm{BB}(x)$ for \cref{sec:flow:split} onward. Using \texttt{out}
instead would silently substitute a lossy, non-invertible quantity for the one the theorem requires.

As a check: $3072\times7\times7=150528=3\times224\times224$, matching the raw pixel count exactly and
confirming dimension-preservation end-to-end, consistent with every row of the table above.

\paragraph{Inverse.} \texttt{iRevNet.inverse} runs the $100$ blocks in reverse order via
\texttt{irevnet\_block.inverse} (subtract $F$; undo \texttt{psi} with \texttt{psi.inverse}, only on the
stride-$2$ blocks), then \texttt{init\_psi.inverse} last -- mirroring Steps~$2\to0$ exactly.

\subsection{The Jacobian of Backbone}
\label{app:irevnet:jacobian}
\textbf{Claim: $\log|\det J_{\mathrm{BB}}(x)|=0$ exactly, for every $x$}, derived from the code above (under
the bijective, \texttt{injective\_pad}-free config of \cref{app:irevnet:checkpoint}), not merely asserted
from the additive-coupling architecture family in the abstract.

\paragraph{Step 1: every \texttt{psi} call has $|\det|=1$.} \texttt{psi.forward}/\texttt{.inverse}
(\texttt{model\_utils.py}) are pure \texttt{reshape}+\texttt{permute}: every output entry equals exactly one
input entry, relocated, with no arithmetic combination of values. As a linear map on flattened coordinates
this is a permutation matrix, so $\det=\pm1$ regardless of input. This covers \texttt{init\_psi} (once, Step
0) and every per-block \texttt{psi} call (once per stage, in that stage's stride-$2$ block, applied to
\emph{both} streams).

\paragraph{Step 2: every block has $|\det|=1$ -- including the part that is not obvious in advance.} The
code's actual output order is \texttt{(x2, y1)} with \texttt{y1 = Fx2 + x1}, i.e.\ the map
$(x_1,x_2)\mapsto(x_2,\ x_1+F(x_2))$, whose Jacobian in block form is
\begin{equation}
J=\begin{pmatrix}0 & I\\ I & \partial F/\partial x_2\end{pmatrix}.
\label{eq:irn-block-jac}
\end{equation}
Swapping the two row-blocks gives $\begin{pmatrix}I & \partial F/\partial x_2\\ 0 & I\end{pmatrix}$ --
upper-triangular, identity diagonal, $\det=1$ -- and a row-block swap only changes the sign of the
determinant, never its magnitude. So $|\det J_{\text{block}}|=1$ exactly, and 
$\partial F/\partial x_2$ -- everything the BatchNorm/ReLU/conv stack inside \linebreak
\texttt{bottleneck\_block} computes -- cancels out of the
determinant entirely; this is also why $F$ itself never needs to be invertible. For a stride-$2$ block, 
this structure composes with the (Step~1) \texttt{psi} application on both streams; a product of two
unit-magnitude determinants is itself unit-magnitude.

\paragraph{Step 3: the final \texttt{merge}.} \texttt{out\_bij = merge(out[0], out[1])} is
\texttt{torch.cat} -- again a pure reshape, $|\det|=1$.

\paragraph{Step 4: compose across the whole network.} Log-determinant of a composition is the sum of the
pieces' log-determinants:
\begin{equation}
\log\bigl|\det J_{\mathrm{BB}}(x)\bigr| = \underbrace{0}_{\texttt{init\_psi}} \!\!+
\sum_{\substack{\text{100}\\\text{blocks}}} \underbrace{0}_{\text{each block}} \!\!+\; 
\underbrace{0}_{\texttt{merge}} \;=\; 0,
\label{eq:irn-total-jac}
\end{equation}
identically in $x$. \Cref{eq:irn-total-jac} is what \cref{eq:flow-jac-split,eq:flow-loss} use as
$\log|\det J_{\mathrm{BB}}(x)|=0$; \cref{rem:irevnet-structure}'s open question about whether the per-
stage $\mathcal S_j$ operator (identified there as \texttt{psi}) might carry its own non-trivial log-det 
term is resolved by Step~1 above: it does not, in this implementation.

\paragraph{Scope of this derivation.} It assumes the \texttt{injective\_pad}-free channel schedule 
confirmed
in \cref{app:irevnet:checkpoint} for the actual released checkpoint. A different \texttt{nChannels}
configuration for which \texttt{pad != 0} at some \texttt{stride=1} block would introduce a genuinely
dimension-expanding padding step there, and \cref{eq:irn-total-jac} would not apply to that variant 
without separately accounting for it.
